\documentclass[10pt,a4paper]{amsart}
\usepackage{amsmath,amssymb,amsthm,graphicx}
\usepackage{verbatim}
\newenvironment{reference}{\comment}{\endcomment}
%\newenvironment{reference}{}{}
\newenvironment{slogan}{\comment}{\endcomment}
\newenvironment{history}{\comment}{\endcomment}
% For commutative diagrams we use Xy-pic
\usepackage[all]{xy}

% We use 2cell for 2-commutative diagrams.
\xyoption{2cell}
% Theorem environments.
%
\theoremstyle{plain}
\newtheorem{theorem}[subsection]{Theorem}
\newtheorem{proposition}[subsection]{Proposition}
\newtheorem{lemma}[subsection]{Lemma}

\theoremstyle{definition}
\newtheorem{definition}[subsection]{Definition}
\newtheorem{example}[subsection]{Example}
\newtheorem{exercise}[subsection]{Exercise}
\newtheorem{situation}[subsection]{Situation}

\theoremstyle{remark}
\newtheorem{remark}[subsection]{Remark}
\newtheorem{remarks}[subsection]{Remarks}

\numberwithin{equation}{subsection}

% Macros
%
\def\lim{\mathop{\mathrm{lim}}\nolimits}
\def\colim{\mathop{\mathrm{colim}}\nolimits}
\def\Spec{\mathop{\mathrm{Spec}}}
\def\Hom{\mathop{\mathrm{Hom}}\nolimits}
\def\Ext{\mathop{\mathrm{Ext}}\nolimits}
\def\SheafHom{\mathop{\mathcal{H}\!\mathit{om}}\nolimits}
\def\SheafExt{\mathop{\mathcal{E}\!\mathit{xt}}\nolimits}
\def\Sch{\mathit{Sch}}
\def\Mor{\mathop{\mathrm{Mor}}\nolimits}
\def\Ob{\mathop{\mathrm{Ob}}\nolimits}
\def\Sh{\mathop{\mathit{Sh}}\nolimits}
\def\NL{\mathop{N\!L}\nolimits}
\def\CH{\mathop{\mathrm{CH}}\nolimits}
\def\proetale{{pro\text{-}\acute{e}tale}}
\def\etale{{\acute{e}tale}}
\def\QCoh{\mathit{QCoh}}
\def\Ker{\mathop{\mathrm{Ker}}}
\def\Im{\mathop{\mathrm{Im}}}
\def\Coker{\mathop{\mathrm{Coker}}}
\def\Coim{\mathop{\mathrm{Coim}}}

% Boxtimes
%
\DeclareMathSymbol{\boxtimes}{\mathbin}{AMSa}{"02}

%
% Macros for moduli stacks/spaces
%
\def\QCohstack{\mathcal{QC}\!\mathit{oh}}
\def\Cohstack{\mathcal{C}\!\mathit{oh}}
\def\Spacesstack{\mathcal{S}\!\mathit{paces}}
\def\Quotfunctor{\mathrm{Quot}}
\def\Hilbfunctor{\mathrm{Hilb}}
\def\Curvesstack{\mathcal{C}\!\mathit{urves}}
\def\Polarizedstack{\mathcal{P}\!\mathit{olarized}}
\def\Complexesstack{\mathcal{C}\!\mathit{omplexes}}
% \Pic is the operator that assigns to X its picard group, usage \Pic(X)
% \Picardstack_{X/B} denotes the Picard stack of X over B
% \Picardfunctor_{X/B} denotes the Picard functor of X over B
\def\Pic{\mathop{\mathrm{Pic}}\nolimits}
\def\Picardstack{\mathcal{P}\!\mathit{ic}}
\def\Picardfunctor{\mathrm{Pic}}
\def\Deformationcategory{\mathcal{D}\!\mathit{ef}}
\usepackage{cleveref}
\crefname{originalreference}{original reference}{original references}
\Crefname{originalreference}{Original reference}{Original references}
\makeatletter
\newlabel{algebraic-lemma-map-space-into-stack}{{algebraic-lemma-map-space-into-stack}{source}{}{sourcepointer-1}{}}
\newlabel{algebraic-lemma-map-space-into-stack@cref}{{[originalreference][0][]algebraic-lemma-map-space-into-stack}{[1][0][]source}{}{}{}}
\makeatother
\crefname{originalreference}{original reference}{original references}
\Crefname{originalreference}{Original reference}{Original references}
\makeatletter
\newlabel{categories-lemma-diagonal-2}{{categories-lemma-diagonal-2}{source}{}{sourcepointer-2}{}}
\newlabel{categories-lemma-diagonal-2@cref}{{[originalreference][0][]categories-lemma-diagonal-2}{[1][0][]source}{}{}{}}
\makeatother
\crefname{originalreference}{original reference}{original references}
\Crefname{originalreference}{Original reference}{Original references}
\makeatletter
\newlabel{criteria-lemma-BG-algebraic}{{criteria-lemma-BG-algebraic}{source}{}{sourcepointer-3}{}}
\newlabel{criteria-lemma-BG-algebraic@cref}{{[originalreference][0][]criteria-lemma-BG-algebraic}{[1][0][]source}{}{}{}}
\makeatother
\crefname{originalreference}{original reference}{original references}
\Crefname{originalreference}{Original reference}{Original references}
\makeatletter
\newlabel{criteria-lemma-flat-quotient-flat-presentation}{{criteria-lemma-flat-quotient-flat-presentation}{source}{}{sourcepointer-4}{}}
\newlabel{criteria-lemma-flat-quotient-flat-presentation@cref}{{[originalreference][0][]criteria-lemma-flat-quotient-flat-presentation}{[1][0][]source}{}{}{}}
\makeatother
\crefname{originaldefinition}{original definition}{original definitions}
\Crefname{originaldefinition}{Original definition}{Original definitions}
\makeatletter
\newlabel{definition-isom}{{definition-isom}{source}{}{sourcepointer-5}{}}
\newlabel{definition-isom@cref}{{[originaldefinition][0][]definition-isom}{[1][0][]source}{}{}{}}
\makeatother
\crefname{originalreference}{original reference}{original references}
\Crefname{originalreference}{Original reference}{Original references}
\makeatletter
\newlabel{stacks-lemma-characterize-essentially-surjective-when-ff}{{stacks-lemma-characterize-essentially-surjective-when-ff}{source}{}{sourcepointer-6}{}}
\newlabel{stacks-lemma-characterize-essentially-surjective-when-ff@cref}{{[originalreference][0][]stacks-lemma-characterize-essentially-surjective-when-ff}{[1][0][]source}{}{}{}}
\makeatother
\crefname{originalreference}{original reference}{original references}
\Crefname{originalreference}{Original reference}{Original references}
\makeatletter
\newlabel{stacks-lemma-when-gerbe}{{stacks-lemma-when-gerbe}{source}{}{sourcepointer-7}{}}
\newlabel{stacks-lemma-when-gerbe@cref}{{[originalreference][0][]stacks-lemma-when-gerbe}{[1][0][]source}{}{}{}}
\makeatother
\title{A gerbe with a section is the quotient of the source by the group of automorphisms of the section}
\author{The Stacks Project Authors}
\date{stacks-project:stacks-morphisms.tex:local-snapshot-20261012; GNU FDL 1.2 or later}
\begin{document}
\maketitle
\noindent\textit{Editorial note.} All delivered mathematical text is the author's own, copied from the pinned Stacks Project chapter source sources/stacks-algebraic-stacks-morphisms/pinned-original-source.tex; no statement or proof step was rewritten, reordered or removed.\par\medskip
\noindent\textit{Editorial note.} Preamble provenance: the Stacks Project keeps one single shared preamble for every chapter; this card takes it directly from the staged sources/stacks-algebraic-stacks-morphisms/preamble.tex as three verbatim spans (lines 8-12, 14-18 and 154-217), so no derived file was written and nothing was added or redefined. The wrapper class and the base packages are supplied by the assembler.\par\medskip
\noindent\textit{Editorial note.} Bibliography provenance: neither the counted statement nor its proof contains a citation (this lane rescans the whole span before writing the card), so the card carries no bibliography and no entry was reconstructed.\par\medskip
\noindent\textit{Editorial note.} Selection provenance: this card was opened under the file-based control contract (CONTROL.md), with the mechanical pre-card gate run on its statement and proof interval before the card was written.\par\medskip
\noindent\textit{Allowed standard prerequisites (editorial).}\begin{itemize}
\item The description of the automorphism group space of an object of a stack as a base change of the inertia. Used inside the counted proof (source lines 6203-6215). Identifies the group acting on the source of the section.
\item The construction of the quotient stack of a group action and the identification of its points. Invoked inside the counted proof (source lines 6216-6230). Produces the quotient stack and the comparison morphism.
\item The two gerbe clauses and the fact that a morphism with a section is surjective. Used at the end of the counted proof (source lines 6231-6244). Shows that the comparison morphism is an equivalence by checking objects and morphisms.
\end{itemize}\medskip
\ifdefined\hypertarget\hypertarget{sourcepointer-1}{}\fi
\noindent\textit{Source reference.} Original reference algebraic-lemma-map-space-into-stack: sources/stacks-algebraic-stacks-morphisms/pinned-original-source.tex, source lines 3831--3831. Navigation pointer only: the counted text refers to the Stacks item algebraic-lemma-map-space-into-stack, which belongs to another chapter of the project and is not part of this card. The pointer carries the target's own label as its number because this chapter's staged snapshot does not contain the other chapter's tag table; its content is not reproduced here.\par\medskip
\ifdefined\hypertarget\hypertarget{sourcepointer-2}{}\fi
\noindent\textit{Source reference.} Original reference categories-lemma-diagonal-2: sources/stacks-algebraic-stacks-morphisms/pinned-original-source.tex, source lines 6218--6218. Navigation pointer only: the counted text refers to the Stacks item categories-lemma-diagonal-2, which belongs to another chapter of the project and is not part of this card. The pointer carries the target's own label as its number because this chapter's staged snapshot does not contain the other chapter's tag table; its content is not reproduced here.\par\medskip
\ifdefined\hypertarget\hypertarget{sourcepointer-3}{}\fi
\noindent\textit{Source reference.} Original reference criteria-lemma-BG-algebraic: sources/stacks-algebraic-stacks-morphisms/pinned-original-source.tex, source lines 6060--6060. Navigation pointer only: the counted text refers to the Stacks item criteria-lemma-BG-algebraic, which belongs to another chapter of the project and is not part of this card. The pointer carries the target's own label as its number because this chapter's staged snapshot does not contain the other chapter's tag table; its content is not reproduced here.\par\medskip
\ifdefined\hypertarget\hypertarget{sourcepointer-4}{}\fi
\noindent\textit{Source reference.} Original reference criteria-lemma-flat-quotient-flat-presentation: sources/stacks-algebraic-stacks-morphisms/pinned-original-source.tex, source lines 3877--3877. Navigation pointer only: the counted text refers to the Stacks item criteria-lemma-flat-quotient-flat-presentation, which belongs to another chapter of the project and is not part of this card. The pointer carries the target's own label as its number because this chapter's staged snapshot does not contain the other chapter's tag table; its content is not reproduced here.\par\medskip
\ifdefined\hypertarget\hypertarget{sourcepointer-5}{}\fi
\noindent\textit{Source reference.} Original definition definition-isom: sources/stacks-algebraic-stacks-morphisms/pinned-original-source.tex, source lines 1357--1357. Navigation pointer only: the counted text refers to the Stacks item definition-isom, which is an existing labeled result of this same chapter and is not part of this card. The pointer carries the target's own label as its number and its definition line as its span; its content is not reproduced here.\par\medskip
\ifdefined\hypertarget\hypertarget{sourcepointer-6}{}\fi
\noindent\textit{Source reference.} Original reference stacks-lemma-characterize-essentially-surjective-when-ff: sources/stacks-algebraic-stacks-morphisms/pinned-original-source.tex, source lines 6232--6232. Navigation pointer only: the counted text refers to the Stacks item stacks-lemma-characterize-essentially-surjective-when-ff, which belongs to another chapter of the project and is not part of this card. The pointer carries the target's own label as its number because this chapter's staged snapshot does not contain the other chapter's tag table; its content is not reproduced here.\par\medskip
\ifdefined\hypertarget\hypertarget{sourcepointer-7}{}\fi
\noindent\textit{Source reference.} Original reference stacks-lemma-when-gerbe: sources/stacks-algebraic-stacks-morphisms/pinned-original-source.tex, source lines 6111--6111. Navigation pointer only: the counted text refers to the Stacks item stacks-lemma-when-gerbe, which belongs to another chapter of the project and is not part of this card. The pointer carries the target's own label as its number because this chapter's staged snapshot does not contain the other chapter's tag table; its content is not reproduced here.\par\medskip
\begin{lemma}
\label{lemma-gerbe-with-section}
Let $\pi : \mathcal{X} \to U$ be a morphism from an algebraic stack to
an algebraic space and let $x : U \to \mathcal{X}$ be a section of $\pi$.
Set $G = \mathit{Isom}_\mathcal{X}(x, x)$, see
Definition \ref{definition-isom}.
If $\mathcal{X}$ is a gerbe over $U$, then
\begin{enumerate}
\item there is a canonical equivalence of stacks in groupoids
$$
x_{can} : [U/G] \longrightarrow \mathcal{X}.
$$
where $[U/G]$ is the quotient stack for the trivial
action of $G$ on $U$,
\item $G \to U$ is flat and locally of finite presentation, and
\item $U \to \mathcal{X}$ is surjective, flat, and locally of finite
presentation.
\end{enumerate}
\end{lemma}

\begin{proof}
Set $R = U \times_{x, \mathcal{X}, x} U$. The morphism $R \to U \times U$
factors through the diagonal $\Delta_U : U \to U \times U$ as it factors
through $U \times_U U = U$. Hence $R = G$ because
\begin{align*}
G & = \mathit{Isom}_\mathcal{X}(x, x) \\
& = U \times_{x, \mathcal{X}} \mathcal{I}_\mathcal{X} \\
& = U \times_{x, \mathcal{X}}
(\mathcal{X}
\times_{\Delta, \mathcal{X} \times_S \mathcal{X}, \Delta}
\mathcal{X}) \\
& = (U \times_{x, \mathcal{X}, x} U) \times_{U \times U, \Delta_U} U \\
& = R \times_{U \times U, \Delta_U} U \\
& = R
\end{align*}
for the fourth equality use
Categories, Lemma \ref{categories-lemma-diagonal-2}.
Let $t, s : R \to U$ be the projections.
The composition law $c : R \times_{s, U, t} R \to R$ constructed on $R$ in
Algebraic Stacks, Lemma \ref{algebraic-lemma-map-space-into-stack}
agrees with the group law on $G$ (proof omitted). Thus
Algebraic Stacks, Lemma \ref{algebraic-lemma-map-space-into-stack}
shows we obtain a canonical fully faithful $1$-morphism
$$
x_{can} : [U/G] \longrightarrow \mathcal{X}
$$
of stacks in groupoids over $(\Sch/S)_{fppf}$. To see that it is
an equivalence it suffices to show that it is essentially surjective.
To do this it suffices to show that any object of $\mathcal{X}$ over
a scheme $T$ comes fppf locally from $x$ via a morphism $T \to U$, see
Stacks, Lemma \ref{stacks-lemma-characterize-essentially-surjective-when-ff}.
However, this follows the condition that $\pi$ turns $\mathcal{X}$
into a gerbe over $U$, see property (2)(a) of
Stacks, Lemma \ref{stacks-lemma-when-gerbe}.

\medskip\noindent
By
Criteria for Representability, Lemma \ref{criteria-lemma-BG-algebraic}
we conclude that $G \to U$ is flat and locally of finite presentation.
Finally, $U \to \mathcal{X}$ is surjective, flat, and locally of finite
presentation by
Criteria for Representability, Lemma
\ref{criteria-lemma-flat-quotient-flat-presentation}.
\end{proof}

\end{document}
